\guidechapter{Start with the problem, not the symbols}{First pass: ordinary arithmetic is enough for this opening.}

You have two long strings of binary digits. You want their exact product.
The research question is how the required work grows when the strings become
longer. It is not a question about whether multiplication is possible, and it
is not a promise that a new program will beat your calculator tomorrow.

Think of three bills. There is a bill for combining numbers, a bill for moving
the right pieces next to one another, and a bill for making sure the final
answer is exact. This paper proposes a way to reduce the first two bills
together while still paying the third. The illustrations below show actual
data paths: the left picture changes where records are stored; the right
picture connects pairs that are combined at successive transform levels.

\explfigure{11-proof-roadmap}{Two distinct jobs inside fast multiplication.}{
Left: records must be brought into a useful order. Right: an eight-entry
butterfly network combines entries differing in one address bit per level.
These are small structural illustrations, not the paper's enormous circuit.}

\paragraph{What is a bit?}
A bit is a zero or a one. In decimal, the places are ones, tens, hundreds.
In binary, they are ones, twos, fours, eights. Thus binary $1011$ means
$1\cdot8+0\cdot4+1\cdot2+1\cdot1=11$. This input has four bits.
The paper writes $n$ for the bit length, not the number's value. Leading
zeroes are allowed: $0011$ is also a four-bit input, even though its value is three.

\paragraph{What does a logarithm count?}
For powers of two, $\log_2 n$ counts how many doublings take one to $n$.
Since $1,2,4,8,16$ takes four doublings, $\log_2 16=4$.
The distinction between $n$ and $\log_2 n$ is crucial. When a bit string
doubles in length, $n$ doubles but $\log_2 n$ increases by just one.

\paragraph{What does a running-time formula mean?}
Schoolbook multiplication pairs each of roughly $n$ digits with roughly
$n$ others, suggesting $n^2$ elementary digit products. More sophisticated
algorithms avoid doing all those pairs separately. A bound such as
$O(n\log n)$ says that, beyond a threshold, the work is at most a fixed
constant times $n\log n$. It does not specify the constant, the threshold,
or the number of seconds on a particular computer.

\explfigure{12-growth-rates}{Growth rates need the right scale.}{
Left: the vertical axis is the logarithm of a formula's value, so exponentially
different values fit on one plot. Right: dividing the toy improved formula
by $n\log_2 n$ exposes its smaller growth. The saving $1/4$ is deliberately
large enough to see; the original paper uses $2^{-182}$. Neither panel is a
measurement of a multiplication program.}

\paragraph{Two ways to read this edition.}
For a first pass, follow the worked examples, pictures, and ``What to carry
forward'' paragraphs. They build the vocabulary before the technical argument
uses it. For a second pass, read the equations and then the corresponding
original proofs in the expanded paper. Understanding why the pieces fit is a
different goal from checking every theorem. The guide supports the first goal
and provides a route toward the second; it does not turn the full proof into
elementary arithmetic.

\paragraph{A quick check.}
If an input grows from 16 bits to 32 bits, $n$ doubles and $\log_2 n$ grows
from four to five. The value of the input integer can be vastly larger than 32.
Keeping these three quantities separate prevents many later confusions.

\takeaway{The paper tries to save a little work at an enormous number of
scales. To understand it, follow both the arithmetic and the movement of data.}
